\documentclass[11pt,a4paper]{article}
\usepackage[utf8]{inputenc}
\usepackage[T1]{fontenc}
\usepackage[spanish,es-nodecimaldot]{babel}
\usepackage{geometry,graphicx,booktabs,amsmath,hyperref,xcolor,listings}
\geometry{margin=2.5cm}
\hypersetup{colorlinks=true,urlcolor=blue,linkcolor=black}
\lstset{basicstyle=\ttfamily\small,frame=single,breaklines=true}

\begin{document}

%% ============ CARATULA ============
\begin{titlepage}
\centering
% Pega el logo como logo_utec.png en esta carpeta; mientras tanto se muestra un recuadro.
\IfFileExists{logo_utec.png}{\includegraphics[width=5cm]{logo_utec.png}}{\fbox{\parbox[c][3cm][c]{5cm}{\centering LOGO UTEC\\(logo\_utec.png)}}}\\[1.5cm]
{\large UNIVERSIDAD DE INGENIERÍA Y TECNOLOGÍA}\\[0.3cm]
{\large Departamento de Ciencia de la Computación}\\[1.5cm]
{\Large \textbf{CS2023 -- Algoritmos y Estructuras de Datos}}\\[0.3cm]
Periodo 2026-2\\[1.5cm]
{\LARGE \textbf{Proyecto Final 2: Anima tu Algoritmo o Estructura de Datos}}\\[0.8cm]
{\Large Estructura asignada: \textbf{Puentes} (algoritmo de Tarjan con \texttt{tin}/\texttt{low})}\\[1.5cm]
\textbf{Integrantes}\\[0.3cm]
\begin{tabular}{ll}
{[Nombre y apellidos Integrante A]} & {[código]}\\
{[Nombre y apellidos Integrante B]} & {[código]}\\
{[Nombre y apellidos Integrante C, si aplica]} & {[código]}\\
\end{tabular}\\[1.5cm]
Docente: Wilder Nina Choquehuayta\\[0.3cm]
Lima, noviembre de 2026
\end{titlepage}

%% ============ PARTE B ============
% Las secciones de descripción del TDA/algoritmo, contribuciones y conclusiones las redactan A y/o C.

\section{Casos borde y comparación (Parte B)}
Esta parte del proyecto valida la implementación de Tarjan en situaciones límite y la compara con una
fuerza bruta real. Ambas implementaciones están escritas por el grupo en C++ y emiten una traza JSON
con el mismo formato; las escenas de Manim solo reproducen esa traza, sin calcular nada por su cuenta.

\subsection{Fuerza bruta}
Para cada arista $e$ del grafo se cuenta el número de componentes conexas $c_0$ del grafo original y el
número $c_e$ del grafo sin $e$, usando BFS lanzado desde cada nodo no visitado. La arista $e$ es un puente
si y solo si $c_e > c_0$. La arista se omite por su identificador, de modo que las aristas paralelas
se tratan correctamente, y el recorrido desde cada nodo sin visitar cubre los grafos desconectados.

\subsection{Casos borde}
\begin{center}
\begin{tabular}{lll}
\toprule
Caso & Resultado esperado & Motivo\\
\midrule
Árbol ($V=6$, $E=5$) & 5 puentes & ninguna back edge; todas las aristas son únicas\\
Ciclo ($V=6$, $E=6$) & 0 puentes & cada arista tiene un camino alterno\\
Aristas paralelas & las paralelas no son puentes & quitar una copia no desconecta\\
Grafo desconectado & puentes solo dentro de cada componente & un DFS por componente\\
\bottomrule
\end{tabular}
\end{center}
En los cuatro casos los puentes de Tarjan coinciden con los de la fuerza bruta. Además,
\texttt{scripts/cross\_check.py} compara ambas implementaciones en 500 grafos aleatorios (con aristas
paralelas, autolazos y nodos aislados) y no encontró diferencias.

\section{Análisis de complejidad}
Sea $V$ el número de vértices y $E$ el de aristas.

\subsection{Tarjan}
Cada vértice se descubre una sola vez y cada arista se examina dos veces (una desde cada extremo) en
las listas de adyacencia. Las actualizaciones de \texttt{tin} y \texttt{low} y la prueba
$\texttt{low}[v] > \texttt{tin}[u]$ cuestan $O(1)$. Con el DFS lanzado desde cada nodo no visitado,
\[ T_{\text{Tarjan}}(V,E) = O(V+E). \]
La memoria es $O(V+E)$ por el grafo, más $O(V)$ para \texttt{tin}, \texttt{low} y la pila de recursión.

\subsection{Fuerza bruta}
Se ejecuta un BFS completo, de costo $O(V+E)$, por cada una de las $E$ aristas, además del BFS inicial:
\[ T_{\text{fuerza bruta}}(V,E) = O\big(E\,(V+E)\big). \]
En un grafo conexo y disperso ($E=\Theta(V)$) esto es $\Theta(V^2)$, mientras que Tarjan es $\Theta(V)$.

\subsection{Medición}
Se define una operación como una entrada de adyacencia examinada o un nodo descubierto. Se usaron
grafos aleatorios conexos con $E\approx 1{,}5\,V$ (semilla fija en \texttt{make\_scaling.py}).
\begin{center}
\begin{tabular}{rrrrr}
\toprule
$V$ & $E$ & Tarjan & Fuerza bruta & Razón\\
\midrule
8 & 11 & 30 & 360 & 12,0\\
16 & 23 & 62 & 1\,488 & 24,0\\
32 & 47 & 126 & 6\,048 & 48,0\\
64 & 95 & 253 & 24\,288 & 96,0\\
128 & 191 & 510 & 97\,920 & 192,0\\
256 & 383 & 1\,022 & 392\,448 & 384,0\\
512 & 767 & 2\,045 & 1\,570\,560 & 768,0\\
\bottomrule
\end{tabular}
\end{center}
Al duplicar $V$, las operaciones de Tarjan se duplican (crecimiento lineal) y las de la fuerza bruta se
cuadruplican (crecimiento cuadrático), de modo que la razón entre ambas crece proporcionalmente a $V$.
En el grafo del video ($V=8$, $E=9$) Tarjan realiza 26 operaciones y la fuerza bruta 260.

\section{Herramientas utilizadas}
\begin{itemize}
\item \textbf{C++17} (g++): Tarjan y fuerza bruta, con exportación de la traza en JSON.
\item \textbf{Python 3} y \textbf{Manim Community}: escenas que leen la traza (\texttt{manim/common.py},
\texttt{edge\_cases.py}, \texttt{complexity.py}).
\item \textbf{FFmpeg}: codificación del video en formato MPEG-4.
\item \textbf{Git y GitHub}: control de versiones y entrega del código.
\item \textbf{\LaTeX}: redacción de este informe.
\end{itemize}

\section{Pasos para reproducir el video}
\begin{enumerate}
\item Clonar el repositorio: \texttt{git clone [URL del repositorio]} y entrar a la carpeta.
\item Instalar los requisitos: g++ (C++17), Python 3.9+, FFmpeg y \texttt{pip install manim}.
\item Compilar, generar las trazas, validar y medir:
\begin{lstlisting}
sh scripts/run_all.sh
\end{lstlisting}
\item Renderizar en alta calidad (1080p, \texttt{.mp4}):
\begin{lstlisting}
manim -qh manim/edge_cases.py CasoArbol CasoCiclo CasoParalelas CasoDesconectado
manim -qh manim/complexity.py Complejidad
\end{lstlisting}
Los videos quedan en \texttt{media/videos/}.
\item Para integrar la Parte A, guardar su implementación como \texttt{src/tarjan.cpp}; el script la detecta y la compila automáticamente. Si falta, solo utiliza Tarjan de referencia para pruebas de la Parte B. Después repetir el paso 3.
\end{enumerate}

\section{Evidencia}
\begin{itemize}
\item Repositorio de GitHub: \url{[PEGAR ENLACE]}
\item Video (.mp4): \url{[PEGAR ENLACE: YouTube no listado o Drive]}
\item Salida de la validación cruzada: \texttt{OK: 500 grafos aleatorios, Tarjan y fuerza bruta coinciden.}
\end{itemize}
\begin{center}
% \includegraphics[width=0.8\textwidth]{captura_arbol.png}
\fbox{\parbox[c][4cm][c]{0.8\textwidth}{\centering [Captura del caso borde: árbol / ciclo]}}\\[0.3cm]
\fbox{\parbox[c][4cm][c]{0.8\textwidth}{\centering [Captura del caso borde: aristas paralelas / desconectado]}}\\[0.3cm]
\fbox{\parbox[c][4cm][c]{0.8\textwidth}{\centering [Captura de la gráfica de complejidad]}}
\end{center}

\end{document}
